\scanpage{38}

\deleted{Je dis \unsure{q}} Le torseur en question provient d'un torseur \inserted{$\Pi$} sous \deleted{\unsure{$\mathscr U_A$}} $U_A\subset A^*$, sous-schém. en groupes des sections de norme $1$, en posant

\[
\Pi=\left\{\Theta\in\Gamma\underbrace{SM^{*}}_{\substack{\text{sections}\\\text{de}\\\text{Norme }1\\(N(\Theta)=1)}}\ \middle|\ \begin{aligned}&\Theta\,\deleted{\unsure{u\mapsto u}}\ \text{normalise }A\\&\text{et y induit }u\longmapsto\bar u\end{aligned}\right\}
\]

4.2. Relation avec les symétries.

Soit $(E,\Delta',Q)$, $Q:E\longrightarrow\Delta'$, une structure de similitude str. fidèle. Soit $D$ un sous-fibré en droites de $E$, \deleted{\unsure{tot|tot.}} non isotrope i.e $Q|D$ est non nul sur chaque fibre. Je vais définir un automorphisme

\[
\sigma_D:E\xrightarrow{\sim}E
\]

de $E$ (avec sa structure de similitude), induisant $\mathrm{id}_D$ sur $D$, $-\mathrm{id}$ sur $E/D$. Quitte à tensoriser $E$ par $D^{-1}$, OPS $D$ muni d'une section base $e$, et profiter de celle-ci pour prendre une base $Q(e)$ dans $\Delta'$, donc OPS que $Q$ est \deleted{\unsure{définie par|défini par}} une fonction quadratique scalaire

\[
Q:E\mathbin{-}\underline O
\]

telle que

\[
Q(e)=1,
\]

% archive sheet number: 37
